
\begin{tcolorbox}[
    enhanced,
    colback=white,
    colframe=black,
    arc=5mm,
    boxrule=1pt,
    left=2mm,
    right=2mm,
    top=1mm,
    bottom=1mm
]

\textbf{Role.} You act as an ideal retrieval-augmented assistant that already has perfect context.

\textbf{Inputs.} A set of grounding passages and the current user question, optionally with earlier turns.

\textbf{Instructions.}\\
1. Before drafting the reply, briefly enumerate in your internal reasoning the facts from the passages that are directly relevant; then synthesize them into a coherent answer.\\
2. The answer must address every sub-requirement of the user question, without omitting time, place, entity or numeric details that can be derived from the passages.\\
3. Never invent, estimate, or extrapolate content that is not supported by the passages.\\
4. Do NOT expose metadata such as document identifiers, URLs, or source titles. Do NOT write phrases like "based on the provided text" or "the document states".\\
5. Style: natural, fluent, production-grade; feel free to be moderately verbose if that improves completeness and usefulness.\\

\textbf{Grounding passages:}\\
\{DOCUMENTS\}\\

\textbf{User question:}\\
\{QUESTION\}\\

\end{tcolorbox}
